% !LW recipe = make-deck
% Level 3F (programs track, FPIL): projection, product and union. Theory:
% images, products, unions, and their list equations. Lean: build
% project, product and union.
% Build: make (rules in ../common/deck.mk).
\DocumentMetadata{
  lang        = en-GB,
  pdfstandard = ua-2,
  tagging     = on,
}
\documentclass{ltx-talk}

\input{preamble}

\title{Level 3F: Projection, Product and Union, as Programs}
\author{The Proofs and Programs Club}
\institute{Royal Holloway, University of London}

\begin{document}

\leveltitle{3F}{Projection, Product and Union, as Programs}
  {Programs track: follows FPIL}{After \levellink{2f}}
  {Next: \levellink{4f}}

\levelframe{3F}{Level 3F at a Glance}{%
  \code{project}, \code{product} and \code{union}.}{%
  the sizes of products and unions of the club's tables, by
  computing.}{%
  write functions on pairs, and on lists of any type & \rds{fpil-poly} \\
  define a function by cases, calling itself on the rest of the list &
    \rds{fpil-data} \\
  check a function on examples: \code{\#eval}, and proofs by \code{rfl} &
    \rds{fpil-eval, fpil-props} \\
}

% ---------------------------------------------------------------- theory
\theorypart{Images, products and unions}

\begin{frame}
  \frametitle{Projection Keeps the Columns You Ask For}
  \defn{Definition 3F.1 (Image)}{For a function $f$ on rows, the
    \emph{image} of a table $R$ is $f[R] = \{f(t) \mid t \in R\}$.}
  \defn{Definition 3F.2 (Projection)}{The projections onto the first and
    the second column are the images under $\pi_1(d_1, d_2) = d_1$ and
    $\pi_2(d_1, d_2) = d_2$.}
  {\centering
    \smalltable{$\pi_1[\rel{Attended}]$}{l}{$\att{name}$}
      {Ada \\ Alan \\ Grace}\hspace{16mm}
    \smalltable{$\pi_2[\rel{Attended}]$}{r}{$\att{week}$}
      {1 \\ 2 \\ 3}\par}
  \slideref{\lnk{ahv4}}
  \theorytag{3F}
\end{frame}

\begin{frame}
  \frametitle{An Image Can Shrink a Table, but Never Grow It}
  \defn{Fact 3F.3}{$|f[R]| \leq |R|$.}
  Alan attended twice, but appears once in $\pi_1[\rel{Attended}]$: four
  rows became three names. A list would keep him twice.
  \defn{Fact 3F.4 (Composition)}{$g[f[R]] = (g \circ f)[R]$.}
  \emph{Proof.} For every $s$:
  \[
    \begin{array}{@{}r@{\quad}l@{}}
      s \in g[f[R]] & \iff s = g(u) \text{ for some } u \in f[R] \\
      & \iff s = g(f(t)) \text{ for some } t \in R \iff s \in (g \circ
        f)[R]. \\
    \end{array}
  \]
  \proofend
  \theorytag{3F}
\end{frame}

\begin{frame}
  \frametitle{A Product Pairs Every Row with Every Row}
  \defn{Definition 3F.5 (Product)}{For tables $R$ and $S$:
    \ $R \times S = \{(t, s) \mid t \in R \text{ and } s \in S\}$.}
  Its rows are pairs of rows. For example,
  $((\text{Ada}, \text{Maths}), (1, \text{Lists})) \in \rel{Member} \times
  \rel{Session}$.
  \defn{Fact 3F.6}{$|R \times S| = |R| \cdot |S|$.}
  So $\rel{Member} \times \rel{Session}$ has $3 \cdot 4 = 12$ rows. In SQL:
  \code{FROM Member, Session}.
  \slideref{\lnk{ahv4}}
  \theorytag{3F}
\end{frame}

\begin{frame}
  \frametitle{Union Combines Two Tables; on Lists, It Is Appending}
  \defn{Definition 3F.7 (Union)}{For tables $R$ and $S$ of the same rows:
    \ $R \cup S = \{t \mid t \in R \text{ or } t \in S\}$.}
  \defn{Fact 3F.8}{If $l$ represents $R$ and $m$ represents $S$
    (Definition 1.2), then $l \mathbin{+\!\!+} m$ represents $R \cup S$.}
  \emph{Proof.} The rows of $l \mathbin{+\!\!+} m$ are those of $l$ and
  those of $m$. \proofend
  \par\medskip
  The list $l \mathbin{+\!\!+} m$ may repeat a row, but it still
  represents the right table. In SQL: \code{UNION}.
  \slideref{\lnk{ahv5}}
  \theorytag{3F}
\end{frame}

\begin{frame}
  \frametitle{On Lists, Projection and Product Are Equations Too}
  \defn{Definition 3F.9 (Projection and product on lists)}{For a function
    $f$ on rows:
    \[
      \begin{array}{@{}l@{\;=\;}l@{}}
        \mathit{project}\ f\ [\,] & [\,] \\
        \mathit{project}\ f\ (t :: l) & f(t) :: \mathit{project}\ f\ l \\
        \mathit{product}\ [\,]\ m & [\,] \\
        \mathit{product}\ (t :: l)\ m & \mathit{project}\ (s \mapsto (t,
          s))\ m \mathbin{+\!\!+} \mathit{product}\ l\ m \\
      \end{array}
    \]}
  The last line pairs $t$ with every row of $m$, then does the same for
  the rest of $l$. Level 6 proves that these lists represent $f[R]$ and
  $R \times S$.
  \theorytag{3F}
\end{frame}

% ---------------------------------------------------------------- Lean
\leanpart{Build project, product and union}

\begin{frame*}
  \frametitle{Your Turn: Build project}
  \leanblock{project.lean}
  \begin{itemize}
    \item \ex{3F.1} Replace the \code{sorry}s with Definition 3F.9.
    \item \ex{3F.2} The test prints Alan twice, but the table
      $\pi_1[\rel{Attended}]$ has him once. Why? (Fact 3F.3.)
  \end{itemize}
  \slideref{\lnk{fpil-poly}}
  \leantag{3F}
\end{frame*}

\begin{frame*}
  \frametitle{Your Turn: Build product}
  \leanblock{product.lean}
  \begin{itemize}
    \item \ex{3F.3} Pair \code{a} with each row of \code{s} (use your
      \code{project}), then \code{++} the product of the rest.
    \item \ex{3F.4} Prove the example: 3 members and 4 sessions make 12
      pairs (Fact 3F.6).
  \end{itemize}
  \leantag{3F}
\end{frame*}

\begin{frame*}
  \frametitle{Your Turn: Build union}
  \leanblock{union.lean}
  \begin{itemize}
    \item \ex{3F.5} \code{union} is \code{r ++ s} (Fact 3F.8). Prove the
      example.
    \item \ex{3F.6} The list has 6 rows, but the table it represents has
      3. Why? (Fact 1.3.)
  \end{itemize}
  \leantag{3F}
\end{frame*}

\end{document}
